\documentclass[12pt]{article}
\usepackage[utf8]{inputenc}
\usepackage[italian]{babel}
\author{Gabriele Guido}
\title{Emulatore CUDA}

\begin{document}
\maketitle
\clearpage
\begin{center}
    \section*{Itroduzione}
\end{center}
La programmazione in CUDA \textit{(Compute Unified Device Architecture)} 
permette di eseguire molte operazioni 
in parallelo grazie alla presenza dei numerosi \textit{core} 
all'interno di una GPU \textit{Graphics Processing Unit}. Questi, a differenza dei classici 
\textit{core} di un processore tradizionale, hanno meno operazioni disponibili ma permettono 
l'elaborazione di molti più dati parallelamente. Questo richiede maggiore attenzione 
nella progettazione degli algoritmi: durante l'esecuzione potrebbero esserci delle dipendenze 
dati che, se non trattate correttamente, possono alterare il risultato atteso.
    
\end{document}